% =====================================================================
% APPENDIX B: USE OF AI SCIENTIFIC ASSISTANTS
% =====================================================================
\appendixbanner{APPENDIX B: USE OF AI SCIENTIFIC ASSISTANTS}
\addcontentsline{toc}{section}{Appendix B: Use of AI scientific assistants}

\subsection*{Workflow}

Three AI assistants were used for the theoretical part, each with a defined role and a verification
step (the full log is \texttt{docs/AI\_USAGE\_LOG.md} in the repository):
\begin{enumerate}
  \item \textbf{Perplexity} searched the literature with structured prompts that asked, for every
        source, the full reference, the DOI, open access and the exact claim it supports.
  \item \textbf{NotebookLM}, restricted to the downloaded PDFs, produced (a) an inventory of every file
        with the bibliographic data as printed in the document, and (b) for each topic of the report,
        short statements with the file, the page and a verbatim quote of at most 30 words.
  \item \textbf{Claude} decided which sources to keep from the inventory, wrote the text using only the
        quote-supported part of each statement, and marked every claim still lacking a quote with a
        visible \texttt{[verify]} tag.
\end{enumerate}
The author checks each tagged claim against the PDF before the final version.

\subsection*{Errors detected (candidates for the AI Error Bounty)}

\begin{enumerate}
  \item \textbf{Off-topic sources.} For a request on glioma classification, Perplexity returned a
        cardiac epigenetic reprogramming paper (twice) and a study on lipid metabolism in glioblastoma
        cell lines, neither supporting any requested point. Detected by comparing title and topic with
        the request; removed.
  \item \textbf{Wrong title and authors, repeated by two tools.} Perplexity and then NotebookLM gave
        one of the downloaded files as ``Machine Learning and Statistical Evaluation of Glioma
        Datasets'' by Garbulowski and Komorowski; the extraction step later attributed its content to
        Sánchez-Marqués et al., and the DOAJ record confirms the latter title and authors
        \cite{sanchezmarques2024}. Two independent AI tools agreeing is therefore not a verification.
  \item \textbf{Misattribution.} Parsons et al.\ (2008, \emph{Science}) was presented as a TCGA
        publication by both tools; it was removed from the TCGA sources.
  \item \textbf{Outdated notion of time.} NotebookLM flagged a 2025 publication year as ``a future
        publication year'', although the inventory was made in October 2026.
  \item \textbf{Wrong DOI.} NotebookLM gave 10.1093/neuonc/noad108 for the EANO guideline on molecular
        diagnostic tools; a Crossref check of all DOIs (script \texttt{report/check\_references.py})
        showed that it belongs to an unrelated article on CAR-T cells, and the correct DOI
        (noad100) was confirmed on the EANO website \cite{sahm2023}. The same check corrected the title
        that NotebookLM gave for \cite{varma2006}.
  \item \textbf{One paper, two identities.} Two uploaded files shared the same DOI; the inventory gave
        one of them a different title and garbled authors, so it was discarded as unreliable, while
        the other was correctly identified later \cite{leiriademendonca2025}.
  \item \textbf{Statements broader than their quotes.} Several extracted statements contained numbers
        or claims absent from the supporting quote (for example, glioblastoma incidence and survival
        figures attributed to a file whose quote only gave lower-grade glioma survival). Only the
        quoted part was used; the rest is marked \texttt{[verify]}.
\end{enumerate}

\subsection*{Reflection}

The verification loop changed how the sources were used: the bibliography records, for every entry,
how its metadata were checked, and no claim enters the text on the authority of an assistant alone.
